\documentclass{article} % This command is used to set the type of
% document you are working on such as an article, book, or presenation

\usepackage{geometry} % This package allows the editing of the page layout
\usepackage{amssymb}
\usepackage{amsmath}  % This package allows the use of a large range
% of mathematical formula, commands, and symbols
\usepackage{enumitem}
\usepackage{graphicx}  % This package allows the importing of images

\newcommand{\question}[2][]{
  \begin{flushleft}
    \textbf{Question #1}: \textit{#2}

\end{flushleft}}
\newcommand{\sol}{\textbf{Solution}:} %Use if you want a boldface solution line
\newcommand{\maketitletwo}[2][]{
  \begin{center}
    \Large{\textbf{Assignment #1}

    Topology} % Name of course here
    \vspace{5pt}

    \normalsize{N. Ohayon Rozanes  % Your name here

    \today}        % Change to due date if preferred
    \vspace{15pt}

\end{center}}
\begin{document}
\maketitletwo[3]  % Optional argument is assignment number
%Keep a blank space between maketitletwo and \question[1]

\question[1]{}

\begin{enumerate}[label=(\alph*)]
  \item $\mathcal{T}_1$ and $\mathcal{T}_2$ are topologies on a set $X$. When is
    $\mathrm{id}(x) = x$ a continuous map from $(X, \mathcal{T}_1)$
    to $(X, \mathcal{T}_2)$?

    A map is continuous iff $\forall V \in \mathcal{T}_2$, $\exists U
    \in \mathcal{T}_1$ such that
    \[
      \mathrm{id}^{-1}(V) = U
    \]
    Since $\mathrm{id}(x) = x \Rightarrow \mathrm{id}^{-1}(V) = \{x
    \in X \mid x \in V\} = V$,
    \[
      V = U
    \]
    therefore, the identity map is continuous only when for every open set in
    $\mathcal{T}_2$, the same set is open in $\mathcal{T}_1$, that is
    \[
      \mathcal{T}_2 \subseteq \mathcal{T}_1
    \]

  \item Want to show $\forall \mathcal{T}$ on $Y$, $Y \subset X$, for which
    $\iota : Y \to X$ is continuous, then $\mathcal{T}_Y \subset \mathcal{T}$.

    Our conditions: $\mathcal{T}$ is a topology on $Y$, and
    \[
      \iota^{-1}(V) \in \mathcal{T}_X \quad \forall V \in \mathcal{T}
    \]
    By definition
    \[
      \mathcal{T}_Y = \{Y \cap U \mid U \in \mathcal{T}_X\}
    \]
    Want to show $\mathcal{T}_Y \subset \mathcal{T}$. $\iota^{-1}$ is
    continuous, therefore for
    all $V \in \mathcal{T}$, $V \in \mathcal{T}_X$ as well. Since
    $U \supset Y \cap U \Rightarrow \mathcal{T}_Y \subset \mathcal{T}$.
\end{enumerate}

\question[2]{}

\begin{enumerate}[label=(\alph*)]
  \item $Y \subset X$, $Y$ is open. Show that $U \subset Y$ is open in
    $\mathcal{T}_Y$ iff $U \in \mathcal{T}$.

    ($\Leftarrow$) Because of closure under finite intersection: $$U
    \in \mathcal{T} \Rightarrow U \cap Y \in
    \mathcal{T}$$ since by definition
    $U \cap Y \in \mathcal{T}_Y$  and for
    $U \subset Y \Rightarrow U \cap Y = U$, then
    $$U \cap Y \in \mathcal{T}_Y \Rightarrow U \in \mathcal{T}_Y$$

    ($\Rightarrow$) $V \in \mathcal{T}_Y \Rightarrow \exists U \in
    \mathcal{T}$ such that
    $U \cap Y = V$, since the intersection of open sets is open it follows that
    $V \in \mathcal{T}$.

  \item
    \begin{enumerate}[label=(\roman*)]
      \item $X \setminus Y \in \mathcal{T}$, $U \subset Y$,
        $Y \setminus U \in \mathcal{T}_Y \Rightarrow X \setminus U
        \in \mathcal{T}$
        \[
          Y \setminus U = V \cap Y \quad \text{for } \exists V \in \mathcal{T}
        \]
        \[
          x \in Y,\ x \notin U \iff x \in V,\ x \in Y
        \]
        \[
          x \notin U \iff x \in V \Rightarrow U = V^c \Rightarrow U =
          X \setminus V
        \]
        \[
          X \setminus U = X \setminus (X \setminus V) = X \cap V = V
          \in \mathcal{T}
        \]

      \item $X \setminus U \in \mathcal{T} \Rightarrow Y \setminus U
        \in \mathcal{T}_Y$

        \[
          x \in Y \cap (X \setminus U) \implies
          x \in Y,\ x \in X,\ x \notin U \implies x \in Y \setminus U
        \]
        Since $X \setminus U \in \mathcal{T}$, we have that $Y \cap
        (X \setminus U) \in \mathcal{T}_Y$ and therefore that
        \[
          Y \cap (X \setminus U) = Y \setminus U \in \mathcal{T}_Y
        \]
    \end{enumerate}
\end{enumerate}

\question[3]{}

$(X_1, d_1)$ and $(X_2, d_2)$ are metric spaces.

\begin{enumerate}[label=(\alph*)]
  \item Define a function $d : X_1 \times X_2 \to \mathbb{R}$,
    \[
      d((x_1, x_2), (y_1, y_2)) = \max(d_1(x_1, y_1), d_2(x_2, y_2))
    \]
  \begin{enumerate}[label=\roman*)]
    \item Since $d \geq d_1$ and $d \geq d_2$, and $d_1, d_2 \geq 0$, then
      it follows $d \geq 0$.
    \item $d_1(x_1, y_1) = 0 \Rightarrow x_1 = y_1$ (same with $d_2$), then
      \[
        d((x_1, x_2), (y_1, y_2)) = 0 \Rightarrow d_1 = 0,\ d_2 = 0
        \Rightarrow (x_1, x_2) = (y_1, y_2)
      \]
    \item
      \begin{align*}
        d((y_1, y_2), (x_1, x_2)) &= \max(d_1(y_1, x_1), d_2(y_2, x_2)) \\
        &= \max(d_2(x_2, y_2), d_1(x_1, y_1)) \\
        &= d((x_1, x_2), (y_1, y_2))
      \end{align*}
    \item $(z_1, z_2) \in X_1 \times X_2$. For each $i \in \{1, 2\}$, by
      the triangle inequality for $d_i$,
      \begin{align*}
        d_i(x_i, y_i) &\leq d_i(x_i, z_i) + d_i(z_i, y_i) \\
        &\leq \max(d_1(x_1, z_1), d_2(x_2, z_2))
        + \max(d_1(z_1, y_1), d_2(z_2, y_2)) \\
        &= d((x_1, x_2), (z_1, z_2)) + d((z_1, z_2), (y_1, y_2))
      \end{align*}
      Since this bound holds for both $i = 1$ and $i = 2$, it also holds
      for the larger of the two, so
      \[
        d((x_1, x_2), (y_1, y_2)) = \max(d_1(x_1, y_1), d_2(x_2, y_2))
        \leq d((x_1, x_2), (z_1, z_2)) + d((z_1, z_2), (y_1, y_2))
      \]
  \end{enumerate}

\item The metric topology on $X_1 \times X_2$ is generated by the basis
  \[
    \mathcal{B} = \{B_d((x_1, x_2), r) \mid (x_1, x_2) \in X_1 \times
    X_2,\ r > 0\}
  \]
  For each $i$, the metric topology on $X_i$ is generated by the basis
  \[
    \mathcal{B}_i = \{B_{d_i}(x, r) \mid x \in X_i,\ r > 0\}
  \]
  and the product topology is generated by the basis of products of
  these balls,
  \[
    \mathcal{B}_p = \{U_1 \times U_2 \mid U_1 \in \mathcal{B}_1,\ U_2
    \in \mathcal{B}_2\}
  \]
  that is,
  \[
    \mathcal{B}_p = \{B_{d_1}(x_1, r_1) \times B_{d_2}(x_2, r_2) \mid
    x_1 \in X_1,\ x_2 \in X_2,\ r_1, r_2 > 0\}
  \]
  Want to show that $\mathcal{T}_{\mathcal{B}} = \mathcal{T}_{\mathcal{B}_p}$.

  ($\Rightarrow$) $\mathcal{T}_{\mathcal{B}} \subset
  \mathcal{T}_{\mathcal{B}_p}$
  \begin{align*}
    (y_1, y_2) \in B_d((x_1, x_2), r)
    &\Leftrightarrow r > \max(d_1(x_1, y_1), d_2(x_2, y_2)) \\
    &\Leftrightarrow \underbrace{r > d_1(x_1, y_1)}_{y_1 \in B_{d_1}(x_1, r)}
    \text{ and } \underbrace{r > d_2(x_2, y_2)}_{y_2 \in B_{d_2}(x_2, r)} \\
    &\Leftrightarrow (y_1, y_2) \in B_{d_1}(x_1, r) \times B_{d_2}(x_2, r)
  \end{align*}
  Thus, for any elements in the basis of the metric topology of
  $d$, there exists a basis element in the product topology such that
  $B_d((x_1,x_2), r)= B_{d_1}(x_1, r)\times B_{d_2}(x_2, r)$,
  therefore we have that $\mathcal{B} \subset \mathcal{B}_p \implies
  \mathcal{T} \subset \mathcal{T}_{\mathcal{B}_p}$

  ($\Leftarrow$) $\mathcal{T}_{\mathcal{B}_p} \subset
  \mathcal{T}_{\mathcal{B}}$,

  Take any $B_{d_1}(x_1, r_1) \times B_{d_2}(x_2, r_2) \in
  \mathcal{B}_p$, we can construct an open set in
  $\mathcal{T}_\mathcal{B}$ as follows, $\forall (y_1, y_2) \in
  B_{d_1}(x_1, r_1) \times B_{d_2}(x_2, r_2)$, define $r =
\min(r_1 -d_1(y_1, x_1)) , r_2 - d_2(y_2, x_2))$ take the open set
$U_{y_1, y_2} \in \mathcal{B}$ as $U_{y_1, y_2} = B_d((y_1,
y_2), r)$. Clearly, $U_{y_1, y_2} \subset B_{d_1}(x_1, r_1) \times B_{d_2}(x_2,
r_2)$ since, by the triangle inequalty $\forall (z_1, z_2)\in U_{y_1,
y_2}, d(z_i, x_i) \le d_i(z_i, y_i) + d_i(y_i, x_i) < r_i$ and
therefore, $(z_1, z_2) \in B_{d_1}(x_1, r_1) \times B_{d_2}(x_2,
r_2)$. Then take
\[
  \bigcup_{(y_1, y_2) \in B_{d_1}(x_1, r_1) \times B_{d_2}(x_2, r_2)}
  U_{y_1, y_2}  = B_{d_1}(x_1, r_1) \times B_{d_2}(x_2, r_2)
\]

Since each $U_{y_1, y_2} \in \mathcal{T}_{\mathcal{B}}$ and the
arbitrary union of open sets is open, we have that $B_{d_1}(x_1,
r_1)\times B_{d_2}(x_2, r_2) \in \mathcal{T}_{\mathcal{B}}$
\end{enumerate}

\end{document}
